% Source: https://artofproblemsolving.com/community/c3h1237575_casework
% Stable identity: aops-topic-1237575-q03
Count by the vertex $P$ of the right angle and the directions of the legs. For each lattice point $P$ and each pair of perpendicular directions given by primitive vectors $(p,q)$ and $(-q,p)$, the legs may have any lengths (multiples of the primitive vectors) that keep both other vertices in the grid; each triangle is counted once (the right angle is at a unique vertex, and we fix the orientation by choosing $(p,q)$ with $p>0$, $q\ge0$).
\begin{itemize}
\item Axis-parallel legs ($(1,0)$ and $(0,1)$): for each $P=(x,y)$ the horizontal leg has $3$ choices of the other endpoint in its row and the vertical leg $3$ in its column: $16\cdot3\cdot3=144$.
\item Legs along $(1,1)$ and $(-1,1)$: for each $P$, multiply the numbers of other grid points on the two diagonals through $P$ and add over all $16$ points; this gives $40$ (each corner contributes $3\cdot0=0$, each of the eight edge points $2\cdot1=2$, and each of the four inner points $3\cdot2=6$, for $16+24=40$).
\item Legs along $(1,2)$ and $(-2,1)$, or $(2,1)$ and $(-1,2)$ (each leg of length $\sqrt5$, the only multiple that fits): $16$.
\end{itemize}
Other primitive directions such as $(1,3)$ do not fit both legs in the grid. The total is $144+40+16=\boxed{200}$.
\begin{asy}
usepackage("amsmath");
size(3cm);
defaultpen(fontsize(10pt));
for(int i=0;i<4;++i) for(int j=0;j<4;++j) dot((i,j));
draw((0,0)--(3,0)--(0,2)--cycle,blue); draw((1,1)--(3,3)--(0,2)--cycle,red);
label("\mbox{}",(min(currentpicture,user=true)+max(currentpicture,user=true))/2,black); // keeps the text after the figure black when LaTeX typesets the labels inline
\end{asy}
